\documentclass[12pt]{article}
\usepackage[utf8]{inputenc}
\usepackage[T1]{fontenc}
\usepackage[french]{babel}
\usepackage{geometry}
\usepackage{enumitem}
\usepackage{titlesec}
\usepackage{hyperref}
\geometry{a4paper, margin=2cm}

\title{Réseau de neurones pour le fanorontelo}
\author{Projet IA - Jeu traditionnel malgache}
\date{}

\begin{document}

\maketitle

\section*{Objectif du projet}
Remplacer un algorithme Minimax par un réseau de neurones capable d'apprendre à jouer au fanorontelo.

\section*{1. Qu'est-ce que le fanorontelo ?}
Jeu traditionnel malgache sur un plateau 3x3. Deux joueurs déplacent leurs pions pour aligner 3 pions, mais \textbf{seulement après que tous leurs pions aient bougé au moins une fois}.

\section*{2. Représentation de l'état}
Chaque case : -2, -1, 0, 1, 2 → encodé en one-hot (5 dimensions). \\
Total : $9 \times 5 + 1 = 46$ dimensions (ajout du joueur actuel).

\section*{3. Architecture du réseau}
\item Entrée : 46
  \item Couches : [128, ReLU], [128, ReLU], [64, ReLU]
  \item Deux têtes : départ (9) et arrivée (9), avec softmax

\section*{4. Sélection du coup}
\begin{enumerate}
  \item Calculer $S(d,a) = p_d[d] \times p_a[a]$
  \item Masquer les coups invalides (case occupée, pas de connexion, mauvais joueur)
  \item Choisir $\arg\max S(d,a)$ parmi les valides
\end{enumerate}

\section*{5. Fonction de perte}
\[
\mathcal{L} = -\log(p_d[d^*]) - \log(p_a[a^*])
\]

\section*{6. Apprentissage}
Backpropagation met à jour :
\item Les poids des têtes indépendamment
  \item Les couches partagées avec le gradient cumulé

\section*{7. Étapes de création}
\begin{enumerate}
  \item Générer données avec Minimax
  \item Encoder états
  \item Construire réseau
  \item Entraîner
  \item Tester
\end{enumerate}

\section*{8. Perspectives}
Agent autonome, application, IA culturelle.

\end{document}
